\section{The remaining cubic core}
\label{sec:remaining-cubic-core}

Combining the five-cut theorem with the six-cycle-shore theorem gives the following reduction.
It does not solve arbitrary six-edge cuts.

\begin{corollary}
There is an absolute constant $K$ such that every connected loopless cubic multigraph of rank $r$ that has a cyclic cut of size at most five or has girth at most six satisfies
\[
 c(G)\le2^{r-1}+Kr2^{r/2}.
\]
Consequently the remaining unrestricted sharp-order question can be restricted to cyclically six-edge-connected cubic kernels of girth at least seven.
\end{corollary}

\begin{proof}
The cyclic-cut case is already proved.
Assume there is no cyclic cut of size at most five, and that $r$ is sufficiently large.
If deleting the vertices of a cycle of length $s$ leaves a forest with $q\ge1$ components, the cubic degree count gives
\[
 r=2s-e(G[V(C)])-q+1\le s.
\]
Thus the complement of a cycle of length at most six contains a cycle when $r>6$.
There can be no cycle of length at most five, since its vertex boundary would be a smaller cyclic cut.

A six-cycle has no chord: a chord would give a cycle of length at most four.
Its boundary therefore has size six.
Its complement is connected, since any cyclic component requires all six boundary edges, leaving no edge to another component.
No complementary vertex has two boundary incidences, since moving such a leaf vertex to the six-cycle shore would give a smaller cyclic cut.
The complement is thus a normalized six-pole.
Apply the six-cycle-shore theorem.
All bounded ranks are absorbed into $K$.
\end{proof}
